% A dataset's lifecycle in its source catalogue, as model/lifecycle.py has it
% and status.yaml records it. Steps that keep the state (a rebuild that finds
% the same bytes, a second link, a recheck, a release step) are left out; the
% steps table in File formats lists them.
\begin{tikzpicture}
\node[actor, minimum width=24mm] (draft) at (0,0) {draft};
\node[actor, minimum width=24mm] (built) at (4.2,0) {built};
\node[actor, minimum width=24mm] (available) at (9.2,0) {available};
\node[actor, minimum width=24mm] (frozen) at (13.4,0) {frozen};
\node[artifact, minimum width=24mm] (withdrawn) at (6.7,-3.6) {withdrawn};
\node[artifact, minimum width=24mm] (purged) at (13.4,-3.6) {purged};
\coordinate (start) at (-3.2,0);
\fill[brand] (start) circle (1.2mm);
\draw[flow] (start) -- node[edgelabel,above]{add} (draft);
\draw[flow] (draft) -- node[edgelabel,above]{build} (built);
\draw[flow] (built) -- node[edgelabel,above,align=center]{upload, link,\\materialize} (available);
\draw[flow] (available) -- node[edgelabel,above]{record} (frozen);
\draw[flow] (available.north) to[out=150, in=30] node[edgelabel,below]{change, revise} (built.north);
\draw[flow] (frozen.north) to[out=160, in=20] node[edgelabel,above]{revise} (built.north west);
\begin{scope}[on background layer]
\node[boundarybox, inner sep=4mm, fit=(draft)(frozen), label={[boundarylabel, anchor=north west]south west:in the catalogue}] (in) {};
\end{scope}
\draw[flow] (in.south -| withdrawn) -- node[edgelabel,right]{remove} (withdrawn);
\draw[flow] (withdrawn) -- node[edgelabel,above]{purge} node[edgelabel,below]{after a release without it} (purged);
\node[note, text width=150mm] at (6.7,-5.4) {\texttt{ethos-data catalog} checks each step against the state in \texttt{status.yaml} and records it there. A rebuild that\\finds other bytes is a change; a new revision is built from new bytes. A release adds a step to every dataset that changed.};
\end{tikzpicture}
